\section{Approved Robustness and Interpretation}

This section records only robustness and interpretation statements authorized by the active theory freeze. It does not extend the exact selective-standardization theorem beyond the quadratic baseline.

\subsection{General concave innovation technology}

Let $g$ be differentiable, increasing, and strictly concave on $[0,E]$. As Proposition~\ref{prop:wedge} states, for any $b_2>b_1$ the private common-layer allocation satisfies
\begin{equation}
 x^F(b_2)\leq x^F(b_1),
\end{equation}
while the coordinated symmetric allocation satisfies
\begin{equation}
 x^S(b_2)\geq x^S(b_1).
\end{equation}
These global order statements include corner solutions. If both compared private optima are interior, the first inequality is strict; if both compared coordinated optima are interior, the second inequality is strict. The interior first-order conditions are
\begin{equation}
 (1-\nu b)g'(x^F)=g'(E-x^F),
 \qquad
 (1+b)g'(x^S)=g'(E-x^S),
\end{equation}
but they are used only to establish strict ordering across two policy values. We do not assert a pointwise sign for $dx^F/db$ or $dx^S/db$ for arbitrary general technologies.

The exact polynomial threshold and global uniqueness in Proposition~\ref{prop:selective} rely on the quadratic baseline. We make no generic claim that sufficiently small $C^2$ perturbations preserve strict concavity or uniqueness of the regulator's reduced objective: induced policy curvature can depend on higher-order features of the innovation technology. The general-technology robustness result is therefore the directional allocation ordering, not a general policy theorem.

\subsection{Incomplete transferability}

Suppose common-layer transferability is $t=\lambda b$ with $\lambda\in(0,1]$. Replacing $b$ by $t$ in the relevant quality spillover gives the interior conditions
\begin{equation}
 (1-\nu t)g'(x^F)=g'(E-x^F),
 \qquad
 (1+t)g'(x^S)=g'(E-x^S).
\end{equation}
Thus the same private-versus-coordinated relative-return wedge is indexed by effective transferability. This reparameterization is the approved robustness statement. We do not infer an additional global policy-uniqueness or persistence theorem from it.

\subsection{Endogenous total R\&D}

The fixed-capacity assumption isolates the composition margin. Consider only an extension in which firms choose common effort $x$ and proprietary effort $z$ and face the same marginal total-capacity cost through a term such as $C(x+z)$. At an interior optimum, that common marginal cost appears in both activity first-order conditions and cancels when the two conditions are compared. The resulting relative condition is
\begin{equation}
 (1-\nu b)g'(x)=g'(z).
 \label{eq:endogenousErelative}
\end{equation}
This condition shows that broader scope still lowers the private relative return to common effort. It does not determine the response of total R\&D, the global response of common effort, or the optimal standardization scope. In particular, we do not claim that a steep capacity cost makes the selective-standardization policy theorem persist by continuity. Those questions require a separate model and are outside the frozen result set.

\subsection{Institutional interpretation}

The variable $b$ is the mandatory scope of a common technological layer or interface. It is not a generic compatibility index and does not mean that all firm knowledge becomes legally disclosed. The reduced-form spillover in \eqref{eq:quality} represents greater cross-product usability of improvements made within the standardized layer.

Two application families illustrate the mechanism without changing the theory. In digital systems, $x$ may represent protocol-, API-, or interface-facing innovation and $z$ proprietary implementation or service innovation. In modular manufacturing, $x$ may represent innovation in standardized external interfaces and $z$ innovation in internal modules or proprietary subsystems. Case evidence from Taiwan's machine-tool industry reports that external interface standards shape product-innovation choices and that innovative effort can concentrate on internal interfaces and modules \citep{ChenLiu2005}. We treat this as suggestive institutional consistency, not causal validation.

The empirical implications should respect the theorem's scope. In the quadratic interior baseline, broader standard scope strictly lowers private common-layer R\&D, and the response is stronger when the competitive penalty $\nu$ is larger. Under a general differentiable increasing strictly concave technology, the corresponding prediction is an order statement: private common-layer R\&D is nonincreasing in scope and may remain flat over ranges where a corner binds. In the quadratic baseline, because $\nu$ rises with downstream substitutability $\rho$, stronger rivalry moves the economy toward the selective-standardization region characterized in Figure~\ref{fig:regime-map}. These predictions concern innovation composition rather than total R\&D expenditure. The model does not imply that the realized transferred improvement $b g(x^F(b))$ must itself rise monotonically with scope.
